% NetSphere_Code_Inventory.tex
% Полный перечень объектов кода (C++ ядро + Python слой) с привязкой к файлам.

\documentclass[14pt,a4paper]{extarticle}
\usepackage[T2A]{fontenc}
\usepackage[utf8]{inputenc}
\usepackage[russian]{babel}
\usepackage{geometry}
\usepackage{hyperref}
\usepackage{enumitem}
\usepackage{xcolor}
\usepackage{setspace}
\usepackage{titlesec}

\geometry{left=25mm,right=15mm,top=20mm,bottom=20mm}
\onehalfspacing
\setlength{\parindent}{1.1cm}

\titleformat{\section}[block]{\normalfont\large\bfseries}{\thesection}{1em}{}
\titleformat{\subsection}[block]{\normalfont\normalsize\bfseries}{\thesubsection}{1em}{}
\titleformat{\subsubsection}[block]{\normalfont\normalsize\bfseries}{\thesubsubsection}{1em}{}

\hypersetup{colorlinks=true,linkcolor=black,urlcolor=blue,citecolor=black}

\begin{document}
\begin{titlepage}
    \begin{center}
        \large
        NetSphere\\[10pt]
        \textbf{Реестр объектов кода (C++ + Python)}\\[12pt]
        {\normalsize Техническая документация для сопровождения и рефакторинга}
    \end{center}
    \vfill
    \begin{flushright}
        \normalsize
        Версия документации: текущее состояние репозитория
    \end{flushright}
\end{titlepage}

\tableofcontents
\newpage

\section{Назначение документа}
Документ предназначен для разработчиков, которым требуется быстро понять интерфейсы и ``точки входа'' для выполнения операций в системе NetSphere:
как создать сеть, как добавлять/удалять домены и устройства, какие параметры нужны и в какой последовательности их вызывать.

\subsection{Ограничения по ``изменению интерфейса от версии к версии''}
В репозитории не выделены явные версии API (нет тегов/версий в коде). Поэтому в документе фиксируется интерфейс \emph{текущего} состояния.
При необходимости можно дополнить раздел, если будут предоставлены предыдущие версии DLL/API (или репозиторий с историей).

\section{Карта модулей (по директориям)}
\begin{itemize}[leftmargin=*]
  \item C++ ядро: \texttt{src/NetSphere/}
  \item Тесты: \texttt{tests/NetSphereTests/}
  \item Python слой:
  \begin{itemize}
    \item GUI: \texttt{python/network\_gui.py}
    \item ctypes-обертка: \texttt{python/corporate\_network.py}
    \item генератор HTML по docstring'ам: \texttt{python/generate\_html\_docs.py}
  \end{itemize}
\end{itemize}

\section{C++ ядро: реестр объектов}
\subsection{Базовые сущности}

\subsubsection{Класс \texttt{NetworkEntity}}
\paragraph{Файл(ы):} \texttt{src/NetSphere/NetworkEntity.h}
\paragraph{Назначение:} абстрактная базовая сущность для доменов и устройств.
\paragraph{Состав:}
protected поле \texttt{std::string id}.
\paragraph{Интерфейс:}
\begin{itemize}
  \item \texttt{NetworkEntity(const std::string\& id)} --- конструктор.
  \item \texttt{virtual \~{}NetworkEntity() = default} --- виртуальный деструктор.
  \item \texttt{const std::string\& getId() const}.
  \item \texttt{virtual void printInfo() const = 0}.
  \item \texttt{virtual std::string getType() const = 0}.
\end{itemize}

\subsubsection{Класс \texttt{Device} (абстрактный)}
\paragraph{Файл(ы):} \texttt{src/NetSphere/Device.h}
\paragraph{Назначение:} базовый класс для устройств сети (хранилища, рабочие станции, принтеры).
\paragraph{Состав:}
protected поле \texttt{std::string macAddress}.
\paragraph{Служебные методы (часть интерфейса для наследников):}
\begin{itemize}
  \item \texttt{static bool isValidMacAddress(const std::string\& mac)}.
  \item \texttt{static void validateId(const std::string\& id)}.
\end{itemize}
\paragraph{Интерфейс:}
\begin{itemize}
  \item \texttt{Device(const std::string\& id, const std::string\& mac)} --- валидирует id и MAC.
  \item \texttt{const std::string\& getMacAddress() const}.
  \item \texttt{virtual std::string getType() const override = 0}.
  \item Переопределение \texttt{printInfo()} реализуется в производных классах.
\end{itemize}

\subsection{Домены}
\subsubsection{Класс \texttt{Domain}}
\paragraph{Файл(ы):} \texttt{src/NetSphere/Domain.h}, \texttt{src/NetSphere/Domain.cpp}
\paragraph{Назначение:} контейнер для устройств и поддоменов; управление доступом через администратор домена.
\paragraph{Состав:}
\begin{itemize}
  \item private поле \texttt{std::string adminId}.
  \item private поле \texttt{std::unordered\_map<std::string, std::shared\_ptr<NetworkEntity>> entities}.
  \item private вспомогательные методы:
  \begin{itemize}
    \item \texttt{checkAdminRights(const std::string\& user) const}.
    \item \texttt{validateEntity(const std::shared\_ptr<NetworkEntity>\& entity) const}.
  \end{itemize}
\end{itemize}
\paragraph{Интерфейс:}
\begin{itemize}
  \item \texttt{Domain(const std::string\& id, const std::string\& admin)} --- валидирует id и admin.
  \item \texttt{void addEntity(std::shared\_ptr<NetworkEntity> entity, const std::string\& user)}.
  \item \texttt{void removeEntity(const std::string\& entityId, const std::string\& user)}.
  \item \texttt{std::shared\_ptr<NetworkEntity> findEntity(const std::string\& entityId) const}.
  \item \texttt{const std::string\& getAdminId() const}.
  \item \texttt{size\_t getEntityCount() const}.
  \item \texttt{const std::unordered\_map<std::string, std::shared\_ptr<NetworkEntity>>\& getAllEntities() const}.
  \item \texttt{void printInfo() const override}.
  \item \texttt{std::string getType() const override} (возвращает \texttt{"Domain"}).
  \item \texttt{void printDetailedInfo() const}.
 \end{itemize}

\subsection{Устройства}
\subsubsection{Класс \texttt{DataStorage}}
\paragraph{Файл(ы):} \texttt{src/NetSphere/DataStorage.h}, \texttt{src/NetSphere/DataStorage.cpp}
\paragraph{Назначение:} модель хранилища данных и доверенных пользователей.
\paragraph{Состав:}
\begin{itemize}
  \item private \texttt{double totalSizeMB}.
  \item private \texttt{double usedSizeMB}.
  \item private \texttt{std::vector<std::string> trustedUsers}.
  \item static \texttt{validateSize(double size)} --- ограничения на размер (0 $<$ size $<=$ 1e6 MB).
\end{itemize}
\paragraph{Интерфейс:}
\begin{itemize}
  \item \texttt{DataStorage(const std::string\& id, const std::string\& mac, double totalSize)}.
  \item \texttt{DataStorage\& operator+=(double additionalSize)}.
  \item \texttt{DataStorage\& operator-=(double sizeToFree)}.
  \item \texttt{DataStorage\& operator=(double newSize)} --- задает \texttt{usedSizeMB}.
  \item Сравнения по \texttt{id}: \texttt{operator<}, \texttt{operator==}, \texttt{operator!=}.
  \item \texttt{void addTrustedUser(const std::string\& user)}.
  \item \texttt{void removeTrustedUser(const std::string\& user)}.
  \item \texttt{bool isUserTrusted(const std::string\& user) const}.
  \item Геттеры: \texttt{getTotalSize()}, \texttt{getUsedSize()}, \texttt{getFreeSize()}, \texttt{getTrustedUsers()}.
  \item \texttt{void printInfo() const override}.
  \item \texttt{std::string getType() const override} (возвращает \texttt{"DataStorage"}).
\end{itemize}

\subsubsection{Класс \texttt{Workstation}}
\paragraph{Файл(ы):} \texttt{src/NetSphere/Workstation.h}, \texttt{src/NetSphere/Workstation.cpp}
\paragraph{Назначение:} модель рабочей станции пользователя.
\paragraph{Состав:}
\begin{itemize}
  \item private \texttt{std::string userId}.
  \item private \texttt{time\_t lastPowerOnTime}.
  \item static \texttt{validateUserId(const std::string\& user)} --- проверка непустоты.
\end{itemize}
\paragraph{Интерфейс:}
\begin{itemize}
  \item \texttt{Workstation(const std::string\& id, const std::string\& mac, const std::string\& user, time\_t powerOnTime)}.
  \item \texttt{const std::string\& getUserId() const}.
  \item \texttt{time\_t getLastPowerOnTime() const}.
  \item \texttt{void updatePowerOnTime(time\_t newTime)}.
  \item \texttt{void printInfo() const override}.
  \item \texttt{std::string getType() const override} (возвращает \texttt{"Workstation"}).
\end{itemize}

\subsubsection{Класс \texttt{Printer}}
\paragraph{Файл(ы):} \texttt{src/NetSphere/Printer.h}, \texttt{src/NetSphere/Printer.cpp}
\paragraph{Назначение:} модель сетевого принтера.
\paragraph{Состав:}
унаследовано от \texttt{Device} (поле \texttt{macAddress} и \texttt{id} в базовом \texttt{NetworkEntity}).
\paragraph{Интерфейс:}
\begin{itemize}
  \item \texttt{Printer(const std::string\& id, const std::string\& mac)}.
  \item \texttt{void printInfo() const override}.
  \item \texttt{std::string getType() const override} (возвращает \texttt{"Printer"}).
\end{itemize}

\subsection{Корпоративная сеть}
\subsubsection{Класс \texttt{CorporateNetwork}}
\paragraph{Файл(ы):} \texttt{src/NetSphere/CorporateNetwork.h}, \texttt{src/NetSphere/CorporateNetwork.cpp}
\paragraph{Назначение:} фасад над всей сетью: корневой домен + глобальный поиск и маршрутизация операций в нужный домен.
\paragraph{Состав:}
\begin{itemize}
  \item private \texttt{std::shared\_ptr<Domain> rootDomain}.
  \item private \texttt{std::unordered\_map<std::string, std::shared\_ptr<NetworkEntity>> allEntities}.
  \item private вспомогательные методы:
  \begin{itemize}
    \item \texttt{collectAllEntities(std::shared\_ptr<Domain> domain)}.
    \item \texttt{findDomainRecursive(std::shared\_ptr<Domain> domain, const std::string\& domainId) const}.
    \item \texttt{findDomainContainingEntity(std::shared\_ptr<Domain> domain, const std::string\& entityId) const}.
  \end{itemize}
\end{itemize}
\paragraph{Интерфейс:}
\begin{itemize}
  \item \texttt{CorporateNetwork(const std::string\& rootAdminId)} --- создаёт \texttt{rootDomain} с \texttt{id = "root\_domain"}.
  \item \texttt{std::shared\_ptr<Domain> getRootDomain() const}.
  \item \texttt{void addEntityToDomain(const std::string\& domainId, std::shared\_ptr<NetworkEntity> entity, const std::string\& user)}.
  \item \texttt{void removeEntity(const std::string\& entityId, const std::string\& user)}.
  \item \texttt{std::shared\_ptr<NetworkEntity> findEntity(const std::string\& entityId) const}.
  \item \texttt{void printDomainInfo(const std::string\& domainId = "") const}.
  \item \texttt{void printNetworkInfo() const}.
\end{itemize}

\subsection{Исключения}
\subsubsection{Система исключений \texttt{NetworkException} и производные}
\paragraph{Файл(ы):} \texttt{src/NetSphere/NetworkExceptions.h}
\paragraph{Назначение:} типизация ошибок доступа и валидации для операций над доменами и устройствами.
\paragraph{Объекты:}
\begin{itemize}
  \item \texttt{NetworkException(std::string message)} : базовый класс, наследник \texttt{std::runtime\_error}.
  \item \texttt{AccessDeniedException(std::string message)} : ``Ошибка доступа: \dots''.
  \item \texttt{ValidationException(std::string message)} : ``Ошибка валидации: \dots''.
  \item \texttt{DeviceOperationException(std::string message)} : ``Ошибка операции с устройством: \dots''.
  \item \texttt{DomainOperationException(std::string message)} : ``Ошибка операции с доменом: \dots''.
\end{itemize}

\section{C++: C API (интеграция с Python)}
\subsection{Экспортируемые объекты}
\paragraph{Файл(ы):} \texttt{src/NetSphere/NetSphereWrapper.h}, \texttt{src/NetSphere/NetSphereWrapper.cpp}
\paragraph{Тип \texttt{enum DeviceType}}
\begin{itemize}
  \item \texttt{DEVICE\_DATASTORAGE = 0}
  \item \texttt{DEVICE\_WORKSTATION = 1}
  \item \texttt{DEVICE\_PRINTER = 2}
\end{itemize}
\paragraph{Функции (extern "C"):}
\begin{itemize}
  \item \texttt{void* create\_network(const char* admin\_id)}
  \item \texttt{void delete\_network(void* network)}
  \item \texttt{const char* get\_network\_info(void* network)}
  \item \texttt{int add\_domain(void* network, const char* domain\_id, const char* admin\_id, const char* user)}
  \item \texttt{int remove\_domain(void* network, const char* domain\_id, const char* user)} --- \emph{присутствует в объявлении, но в текущем \texttt{.cpp} нет реализации}.
  \item \texttt{const char* get\_domain\_info(void* network, const char* domain\_id)} --- \emph{присутствует в объявлении, но в текущем \texttt{.cpp} нет реализации}.
  \item \texttt{int add\_device(void* network, const char* domain\_id, int device\_type, const char* id, const char* mac, const char* user, ...)}
  \item \texttt{int remove\_device(void* network, const char* device\_id, const char* user)}
  \item \texttt{const char* get\_device\_info(void* network, const char* device\_id)} --- в реализации возвращается заглушка.
  \item \texttt{const char* find\_entity(void* network, const char* entity\_id)}
  \item \texttt{void free\_string(const char* str)}
  \item \texttt{const char* get\_last\_error()}
\end{itemize}

\subsection{Семантика varargs у \texttt{add\_device}}
В \texttt{NetSphereWrapper.cpp} реализован выбор дополнительных параметров в зависимости от \texttt{device\_type}.
Последовательность параметров после \texttt{user} задаётся так:
\begin{itemize}
  \item \texttt{DEVICE\_DATASTORAGE (0)}: далее \texttt{double total\_size} (объем MB) .
  \item \texttt{DEVICE\_WORKSTATION (1)}: далее \texttt{const char* user\_id}, затем \texttt{time\_t power\_time}.
  \item \texttt{DEVICE\_PRINTER (2)}: дополнительных параметров не требуется.
\end{itemize}

\subsection{Обработка ошибок}
Внутри \texttt{NetSphereWrapper.cpp} используется шаблонная обертка \texttt{handle\_exception}, которая:
сбрасывает \texttt{last\_error}, перехватывает исключения \texttt{std::exception} и устанавливает текст ошибки.

\section{Python слой: реестр объектов}
\subsection{Модуль \texttt{python/corporate\_network.py}}
\paragraph{Файл(ы):} \texttt{python/corporate\_network.py}
\paragraph{Назначение:} ctypes-обертка, возвращающая Python-объекты \texttt{NetworkEntity} с opaque handle и класс фасада \texttt{CorporateNetwork}.

\subsubsection{Типы данных}
\begin{itemize}
  \item \texttt{class DeviceType(IntEnum)}:
  \begin{itemize}
    \item \texttt{UNKNOWN = 0}
    \item \texttt{DATA\_STORAGE = 1}
    \item \texttt{WORKSTATION = 2}
    \item \texttt{PRINTER = 3}
    \item \texttt{DOMAIN = 4}
  \end{itemize}
  \item \texttt{class DeviceInfo(ctypes.Structure)}:
  поля \texttt{id}, \texttt{mac}, \texttt{type}, \texttt{data}.
  \item \texttt{class CorporateNetworkError(Exception)}.
  \item Функция \texttt{decode\_error\_string(error\_ptr) -> str}.
  \item Функция \texttt{load\_dll(dll\_path: Optional[str])}:
  загружает DLL и настраивает \texttt{argtypes/restype} для ожидаемых C-функций.
 \end{itemize}

\subsubsection{Класс \texttt{NetworkEntity}}
\paragraph{Назначение:} handle-обертка над нативной сущностью (устройство/домен).
\paragraph{Интерфейс:}
\begin{itemize}
  \item \texttt{\_\_init\_\_(handle, dll)}.
  \item \texttt{\_\_del\_\_()} --- вызывает \texttt{dll.device\_destroy(handle)} (ожидаемое имя функции задается в \texttt{load\_dll}).
  \item \texttt{storage\_add\_data(size: float) -> bool}.
  \item \texttt{storage\_free\_data(size: float) -> bool}.
  \item \texttt{storage\_add\_trusted\_user(user: str) -> bool}.
  \item \texttt{get\_info() -> Dict[str, Any]}:
  возвращает \texttt{\{ 'id', 'mac', 'type' \}}.
  \item Все методы при ошибке опираются на \texttt{dll.get\_last\_error()} и переводят его через \texttt{decode\_error\_string}.
\end{itemize}

\subsubsection{Класс \texttt{CorporateNetwork}}
\paragraph{Файл(ы):} \texttt{python/corporate\_network.py}
\paragraph{Назначение:} фасад для создания сети и добавления/удаления сущностей.
\paragraph{Интерфейс:}
\begin{itemize}
  \item \texttt{\_\_init\_\_(root\_admin\_id: str, dll\_path: Optional[str]=None)}:
  грузит DLL (\texttt{load\_dll}) и создает handle через \texttt{dll.network\_create(...)}.
  \item \texttt{\_\_del\_\_()} --- вызывает \texttt{dll.network\_destroy(handle)}.
  \item \texttt{create\_domain(domain\_id: str, admin\_id: str) -> NetworkEntity}.
  \item \texttt{create\_storage(device\_id: str, mac: str, total\_size: float) -> NetworkEntity}.
  \item \texttt{create\_workstation(device\_id: str, mac: str, user\_id: str, power\_on\_time: int) -> NetworkEntity}.
  \item \texttt{create\_printer(device\_id: str, mac: str) -> NetworkEntity}.
  \item \texttt{add\_device\_to\_domain(domain\_id: str, device: NetworkEntity, user: str) -> bool}.
  \item \texttt{get\_last\_error() -> str}.
  \item \texttt{clear\_error()}.
  \item \texttt{remove\_device\_from\_network(device\_id: str, user: str) -> bool}.
\end{itemize}

\subsection{Несовпадение ожидаемого DLL C API}
В \texttt{python/corporate\_network.py} в \texttt{load\_dll} ожидаются экспортируемые функции:
\texttt{network\_create}, \texttt{device\_create\_domain}, \texttt{storage\_add\_data}, \texttt{network\_add\_device} и т.д.

При этом C++ слой \texttt{src/NetSphere/NetSphereWrapper.cpp} экспортирует функции:
\texttt{create\_network}, \texttt{add\_domain}, \texttt{add\_device}, \texttt{remove\_device}, \texttt{get\_device\_info}, \texttt{get\_last\_error} и т.д.

С точки зрения разработчика это означает необходимость:
либо привести Python-ctypes-обертку к именам C API из \texttt{NetSphereWrapper.h},
либо обеспечить в DLL алиасы/обертки с ожидаемыми именами.

\subsection{Модуль \texttt{python/network\_gui.py}}
\paragraph{Файл(ы):} \texttt{python/network\_gui.py}
\paragraph{Назначение:} PyQt5 GUI поверх \texttt{corporate\_network}.
\paragraph{Ключевые объекты:}
\begin{itemize}
  \item \texttt{class DeviceTreeWidget(QTreeWidget)}:
  \begin{itemize}
    \item сигнал \texttt{device\_selected = pyqtSignal(dict)}.
    \item методы: \texttt{\_\_init\_\_}, \texttt{set\_network(network)}, \texttt{on\_item\_clicked(item, column)}.
  \end{itemize}
  \item \texttt{class AddDeviceDialog(QDialog)}:
  \begin{itemize}
    \item методы: \texttt{\_\_init\_\_}, \texttt{setup\_ui()}, \texttt{on\_type\_changed(index)}, \texttt{get\_device\_data()}.
  \end{itemize}
  \item \texttt{class MainWindow(QMainWindow)}:
  \begin{itemize}
    \item методы: \texttt{\_\_init\_\_}, \texttt{init\_ui()}, \texttt{setup\_info\_tab()},
    \texttt{setup\_operations\_tab()}, \texttt{setup\_log\_tab()},
    \texttt{create\_network()}, \texttt{add\_device()}, \texttt{remove\_device()},
    \texttt{on\_device\_selected(device\_info)}, \texttt{on\_add\_data()},
    \texttt{on\_free\_data()}, \texttt{on\_add\_user()}, \texttt{log\_message(message)}.
  \end{itemize}
  \item \texttt{def main()} --- точка входа GUI.
\end{itemize}

\subsection{Модуль \texttt{python/generate\_html\_docs.py}}
\paragraph{Файл(ы):} \texttt{python/generate\_html\_docs.py}
\paragraph{Назначение:} генерация HTML-документации по docstring'ам без импорта модулей.
\paragraph{Ключевые объекты:}
\begin{itemize}
  \item dataclasses: \texttt{FunctionDoc}, \texttt{ClassDoc}, \texttt{ModuleDoc}.
  \item функции:
  \begin{itemize}
    \item \texttt{parse\_module(path) -> ModuleDoc}
    \item \texttt{render\_module\_html(mod, output\_dir)}
    \item \texttt{render\_index\_html(modules, output\_dir)}
    \item \texttt{main()} (CLI).
  \end{itemize}
\end{itemize}

\section{Тесты (NetSphereTests): реестр сценариев}
\paragraph{Файл(ы):} \texttt{tests/NetSphereTests/}
\paragraph{Примечание:} тесты используют GoogleTest и покрывают валидацию и основные сценарии интерфейса.

\subsection{DeviceTests}
\paragraph{Файл:} \texttt{tests/NetSphereTests/DeviceTests.cpp}
\begin{itemize}
  \item \texttt{TEST(DeviceTest, InvalidMacAddress)}
  \item \texttt{TEST(DeviceTest, EmptyMacAddress)}
  \item \texttt{TEST(DeviceTest, InvalidDeviceId)}
  \item \texttt{TEST(DeviceTest, TooLongDeviceId)}
  \item \texttt{TEST(DeviceTest, InvalidCharactersInDeviceId)}
\end{itemize}

\subsection{DataStorageTests}
\paragraph{Файл:} \texttt{tests/NetSphereTests/DataStorageTests.cpp}
\begin{itemize}
  \item \texttt{TEST(DataStorageTest, ConstructorAndGetters)}
  \item \texttt{TEST(DataStorageTest, AdditionOperator)}
  \item \texttt{TEST(DataStorageTest, SubtractionOperator)}
  \item \texttt{TEST(DataStorageTest, AssignmentOperator)}
  \item \texttt{TEST(DataStorageTest, ComparisonOperators)}
  \item \texttt{TEST(DataStorageTest, TrustedUsers)}
  \item \texttt{TEST(DataStorageTest, FreeSizeCalculation)}
  \item \texttt{TEST(DataStorageTest, PrintInfoNoCrash)}
\end{itemize}

\subsection{DataStorageErrorTests}
\paragraph{Файл:} \texttt{tests/NetSphereTests/DataStorageErrorTests.cpp}
\begin{itemize}
  \item \texttt{TEST(DataStorageErrorTest, StorageOverflow)}
  \item \texttt{TEST(DataStorageErrorTest, FreeMoreThanUsed)}
  \item \texttt{TEST(DataStorageErrorTest, NegativeSizeAssignment)}
  \item \texttt{TEST(DataStorageErrorTest, SizeExceedsTotal)}
  \item \texttt{TEST(DataStorageErrorTest, AddEmptyUser)}
  \item \texttt{TEST(DataStorageErrorTest, AddDuplicateUser)}
  \item \texttt{TEST(DataStorageErrorTest, RemoveNonExistentUser)}
  \item \texttt{TEST(DataStorageErrorTest, InvalidStorageSize)}
  \item \texttt{TEST(DataStorageErrorTest, TooLargeStorageSize)}
\end{itemize}

\subsection{DomainTests}
\paragraph{Файл:} \texttt{tests/NetSphereTests/DomainTests.cpp}
\begin{itemize}
  \item \texttt{TEST(DomainTest, ConstructorAndGetters)}
  \item \texttt{TEST(DomainTest, AddAndFindEntities)}
  \item \texttt{TEST(DomainTest, RemoveEntities)}
  \item \texttt{TEST(DomainTest, AddSubdomains)}
  \item \texttt{TEST(DomainTest, GetAllEntities)}
  \item \texttt{TEST(DomainTest, PrintInfoMethods)}
\end{itemize}

\subsection{DomainErrorTests}
\paragraph{Файл:} \texttt{tests/NetSphereTests/DomainErrorTests.cpp}
\begin{itemize}
  \item \texttt{TEST(DomainErrorTest, EmptyDomainId)}
  \item \texttt{TEST(DomainErrorTest, EmptyAdminId)}
  \item \texttt{TEST(DomainErrorTest, AddEntityWithoutAdminRights)}
  \item \texttt{TEST(DomainErrorTest, RemoveEntityWithoutAdminRights)}
  \item \texttt{TEST(DomainErrorTest, AddDuplicateEntity)}
  \item \texttt{TEST(DomainErrorTest, RemoveNonExistentEntity)}
  \item \texttt{TEST(DomainErrorTest, AddNullEntity)}
  \item \texttt{TEST(DomainErrorTest, RemoveEntityWithEmptyId)}
  \item \texttt{TEST(DomainErrorTest, FindEntityWithEmptyId)}
\end{itemize}

\subsection{WorkstationPrinterTests}
\paragraph{Файл:} \texttt{tests/NetSphereTests/WorkstationPrinterTests.cpp}
\begin{itemize}
  \item \texttt{TEST(WorkstationTest, InvalidParameters)}
  \item \texttt{TEST(PrinterTest, InvalidParameters)}
  \item \texttt{TEST(WorkstationTest, WorkstationMethods)}
  \item \texttt{TEST(PrinterTest, PrinterMethods)}
  \item \texttt{TEST(WorkstationTest, UpdatePowerOnTime)}
\end{itemize}

\subsection{CorporateNetworkTests}
\paragraph{Файл:} \texttt{tests/NetSphereTests/CorporateNetworkTests.cpp}
\begin{itemize}
  \item \texttt{TEST(CorporateNetworkTest, ConstructorAndGetters)}
  \item \texttt{TEST(CorporateNetworkTest, AddEntityToRootDomain)}
  \item \texttt{TEST(CorporateNetworkTest, AddEntityToSubdomain)}
  \item \texttt{TEST(CorporateNetworkTest, RemoveEntity)}
  \item \texttt{TEST(CorporateNetworkTest, FindEntity)}
  \item \texttt{TEST(CorporateNetworkTest, PrintDomainInfo)}
  \item \texttt{TEST(CorporateNetworkTest, PrintNetworkInfo)}
  \item \texttt{TEST(CorporateNetworkTest, CollectAllEntities)}
  \item \texttt{TEST(CorporateNetworkTest, ComplexDomainHierarchy)}
\end{itemize}

\subsection{CorporateNetworkErrorTests}
\paragraph{Файл:} \texttt{tests/NetSphereTests/CorporateNetworkErrorTests.cpp}
\begin{itemize}
  \item \texttt{TEST(CorporateNetworkErrorTest, AddEntityToNonExistentDomain)}
  \item \texttt{TEST(CorporateNetworkErrorTest, AddEntityWithoutAccessRights)}
  \item \texttt{TEST(CorporateNetworkErrorTest, RemoveNonExistentEntity)}
  \item \texttt{TEST(CorporateNetworkErrorTest, RemoveEntityWithoutAccessRights)}
  \item \texttt{TEST(CorporateNetworkErrorTest, AddDuplicateEntity)}
  \item \texttt{TEST(CorporateNetworkErrorTest, PrintNonExistentDomainInfo)}
  \item \texttt{TEST(CorporateNetworkErrorTest, AddToSubdomainWithoutAccess)}
  \item \texttt{TEST(CorporateNetworkErrorTest, RemoveFromSubdomainWithoutAccess)}
  \item \texttt{TEST(CorporateNetworkErrorTest, AddNullEntity)}
  \item \texttt{TEST(CorporateNetworkErrorTest, AddEntityWithEmptyId)}
\end{itemize}

\subsection{NetworkExceptionsTests}
\paragraph{Файл:} \texttt{tests/NetSphereTests/NetworkExceptionsTests.cpp}
\begin{itemize}
  \item \texttt{TEST(NetworkExceptionsTest, BaseException)}
  \item \texttt{TEST(NetworkExceptionsTest, AccessDeniedException)}
  \item \texttt{TEST(NetworkExceptionsTest, ValidationException)}
  \item \texttt{TEST(NetworkExceptionsTest, DeviceOperationException)}
  \item \texttt{TEST(NetworkExceptionsTest, DomainOperationException)}
\end{itemize}

\section{Приложение: демонстрационный код (main.cpp)}
\paragraph{Файл:} \texttt{src/NetSphere/main.cpp}
\begin{itemize}
  \item \texttt{void demonstrateDataStorage()}
  \item \texttt{void demonstrateDomainOperations()}
  \item \texttt{void demonstratePolymorphism()}
  \item \texttt{void demonstrateCorporateNetwork()}
  \item \texttt{int main()}
\end{itemize}

\end{document}

